%% start of file `template.tex'.
%% Copyright 2006-2013 Xavier Danaux (xdanaux@gmail.com).
%
% This work may be distributed and/or modified under the
% conditions of the LaTeX Project Public License vergqsion 1.3c,
% available at http://www.latex-project.org/lppl/.


\documentclass[11pt,a4paper,sans]{moderncv}

\moderncvstyle{banking}
\moderncvcolor{blue}

\usepackage{polyglossia}
\usepackage{fontspec}
\usepackage{xunicode}
\usepackage{xltxtra}
\usepackage{amssymb}

\defaultfontfeatures{Scale=MatchLowercase,Mapping=tex-text}
\setdefaultlanguage[spelling=modern]{russian}
\setotherlanguage{english}

\newfontfamily\cyrillicfonttt{Liberation Sans}
\newfontfamily\cyrillicfont{Liberation Sans}

\usepackage[scale=0.77]{geometry}

\name{Алексей}{Уйманов}
\address{Россия}{Тюмень}{}
\phone[mobile]{+7 932-470-79-15}
\email{s9gf4ult@gmail.com}
\jabber{segfault@jabber.ru}
\homepage{https://github.com/s9gf4ult}



\newcommand*{\hlink}[1]{\textcolor{blue}{\texttt{\underline{\href{#1}{#1}}}}}
\newcommand*{\nlink}[2]{\textcolor{blue}{\texttt{\underline{\href{#1}{#2}}}}}
\newcommand*{\bb}[1]{\hbox{#1}}

%% \newcommand*{\pic}[1]{\raisebox{-0.25em}{\makebox[1.5em]{\includegraphics[height=1.1em,keepaspectratio=true]{#1}}}}
\newcommand*{\pic}[1]{\includegraphics[height=1.1em,keepaspectratio=true]{#1}}

\voffset=-0.3in
\textheight=700pt

\begin{document}

\begin{center}
  \namestyle{Алексей~Уйманов}

  \quotestyle{Haskell программист}

  \addressfont\color{color2}%
  Россия, Тюмень
\end{center}

\medskip

\hbox{
\parbox[t]{0.75\linewidth}{
  \begin{itemize}
  \item[\pic{Octocat}] \hlink{https://github.com/s9gf4ult}
  \item[\pic{bucket}] \hlink{https://bitbucket.org/s9gf4ult}
  \item[\pic{haskell}] \hlink{http://hackage.haskell.org/user/AlekseyUymanov}
  \end{itemize}
}%
\parbox[t]{0.25\linewidth}{
  \begin{itemize}
  \item[\mobilephonesymbol] \href{tel:+7-932-470-79-15}{+7 932-470-79-15}
  \item[\emailsymbol] \href{mailto:s9gf4ult@gmail.com}{s9gf4ult@gmail.com}
  \item[\pic{xmpp}] \href{xmpp:segfault@jabber.ru}{segfault@jabber.ru}
  \item[\pic{skype}] \href{skype:s9gf4ult}{s9gf4ult}
  \end{itemize}
}}

\bigskip

\section{Опыт работы}

\cventry{2014 -- Сейчас}{Haskell программист}{AdTel}{Удаленная работа}{}{}{%
  \begin{itemize}
  \item Сделал сервис аутентификации по протоколу \bb{Oauth2}. Сервис
    в последствии был разбит на два пакета: подключаемый сервер-плагин
    \bb{oauth2-server} (\bb{Yesod} subsite) и непосредственно сервис,
    к которому подключен \bb{oauth2-server}.
  \item Сделал пакет для удобной (и безопасной) генерации SQL-запросов
    \nlink{http://hackage.haskell.org/package/postgresql-query}{postgresql-query}. Интерполирующий
    квазиквотер SQL-запросов с ORM-подобной
    функциональностью. Работает через \bb{postgresql-simple}.
  \item Переписал пакет
    \nlink{http://hackage.haskell.org/package/OTP}{OTP} с
    использвоанием \bb{cryptohash}:
    \nlink{http://hackage.haskell.org/package/one-time-password}{one-time-password}.
  \item Сделал пакет для удобного доступа к redis соединению:
    \nlink{http://hackage.haskell.org/package/hedis-monadic}{hedis-monadic}.
  \item Автоматизировал сборку и развертывание сервиса аутентификации
    с использованием
    \nlink{http://hackage.haskell.org/package/shake}{shake}.
  \end{itemize}
}

\cventry{2014}{Ruby on Rails программист}{Студия Кечинова}{Удаленная работа}{}{}{%
  \begin{itemize}
  \item Учавствовал в разработке большого проекта на Ruby on Rails.
  \item Применял git-flow.
  \item Закрыл столько-то задач в трекере ...
  \end{itemize}
}

\cventry{2013 -- 2014}{Программист}{CyberStep}{Тюмень}{}{%
  Разработка программного обеспечения для станков плазменной резки.
  \medskip
  \begin{itemize}
  \item Интерфейс станка плазменной резки (Python, Tcl, C)
    \begin{itemize}
    \item Внедрил базу данных параметров резки.
    \item Сделал модуль интерполяции параметров резки для режимов резки под углом.
    \item Всякая текучка с интерфейсом, исправление багов (Python, Tcl).
    \end{itemize}
  \item Форк Fat Free CRM, добавил генератор комерческих предложений
    для нужд предприятия\\
    \hlink{https://github.com/s9gf4ult/fat\_free\_crm/tree/commoffers}.
  \end{itemize}
}

\cventry{2012 -- 2013}{Веб програмист}{Ирис Интеграция}{Тюмень}{}{%
  \begin{itemize}
  \item Разработал систему электронного оффлайн голосования на Ruby on
    Rails
  \item Адаптировал движок Mediawiki для создания информационного
    портала о городских достопримечательностях
  \end{itemize}
}

\cventry{2009 -- 2012}{Системный администратор}{ГАУ ТО Миац}{Тюмень}{}{%
  Поддержка систем хранения, виртуальных машин.
}

\cventry{2007 -- 2009}{Инженер конструктор}{ОАО Сибнефтемаш}{Тюмень}{}{%
  Прочностной расчет технологических емкостей, конструкций.
}

\cventry{2006 -- 2007}{Системный администратор}{МУПГТ Тюменьгортранс}{Тюмень}{}{}

\newpage

\section{Умения}

\subsection{Программирование}

\begin{itemize}
\item Haskell
  \begin{itemize}
  \item Template Haskell/QuasiQuoter: могу на \hbox{Template Haskell}
    выводить инстансы, генерировать типы, выражения. Делать
    квазиквотеры, наподобие
    \nlink{http://hackage.haskell.org/package/shakespeare}{shakespeare}.
  \item Комбинаторные парсеры: parsec, attoparsec.
  \item Yesod: знаю внутреннее строение на хорошем уровне, делал пулл-реквесты.
  \item Scotty: применял для написания сервиса генерации PDF файлов для AdTel.
  \item QuicCheck, SmallCheck: автоматизированное тестирование свойств
    программы.
  \item MonadBaseControl: знаю зачем нужно и могу написать инстанс
    \includegraphics[height=0.9ex]{Trollface.ps}
  \item Продвинутая магия типов: GADTs, TypeFamilies, RankNTypes.
  \item Профилирование, устранение space leaks.
  \item Имею пакеты на
    \nlink{http://hackage.haskell.org/user/AlekseyUymanov}{хакадже}.
  \item В хаскеле с 2013 года.
  \end{itemize}
\item Web
  \begin{itemize}
  \item CSS: Bootstrap, Blueprint CSS.
  \item JavaScript: JQuery, AngularJS.
  \end{itemize}
\item Python
  \begin{itemize}
  \item PyGtk.
  \item Django.
  \end{itemize}
\item Ruby on Rails
  \begin{itemize}
  \item haml, coffeescript, scss.
  \item rspec, factory\_girl, Faker.
  \item paperclip, simple\_form, devise, acts\_as\_commentable.
  \end{itemize}
\item Низкоуровневые интерфейсы к базам данных: освоил пока писал
  \nlink{http://hackage.haskell.org/package/hdbi}{hdbi}.
\item SQL: на нормальном уровне знаю PostgreSQL и Sqlite, менее хорошо
  знаю MySql.
\item Пишу тесты: unit-тесты, property-тесты, интеграционные.
\item Emacs / ELisp.
\item \LaTeX.
\item Git, Subversion.
\item Redmine, Slack.
\item Имею опыт удаленной работы: самоорганизация / самомотивация.
\end{itemize}

\subsection{Администрирование}

\begin{itemize}
\item Linux.
\item Apache, nginx, Varnish, Squid.
\item Хранение: LVM, RAID, iSCSI.
\item Сеть: iptables, ip, tc.
\item Виртуализация: libvirt, kvm.
\end{itemize}

\section{Образование}
\cventry{2007 -- 2009}{Магистр}{ТюмГНГУ}{Тюмень}{\textit{4.8/5}}{Кафедра прикладной механики.}

\cventry{2003 -- 2007}{Бакалавр}{ТюмГНГУ}{Тюмень}{\textit{4.7/5}}{Кафедра прикладной механики.}

\enlargethispage{2em}

\section{Прочее}
\cvitemwithcomment{Английский}{Beginer}{Свободное чтение
  документации и электронная переписка.}
\cvitemwithcomment{Расположение}{Россия, Тюмень}{Возможность переезда обсуждается.}
\cvitemwithcomment{Возраст}{28}{}

\nocite{*}
\bibliographystyle{plain}
\bibliography{publications}                        % 'publications' is the name of a BibTeX file
\end{document}
